\documentclass{article}
\usepackage[UTF8]{ctex}
\usepackage{amsmath,amssymb,amsfonts}
\usepackage{enumitem}
\usepackage{geometry}
\geometry{a4paper, margin=1in}
% === 错题本排版优化（不参与题目标识） ===
\setlength{\parindent}{0pt}
\setlength{\parskip}{0.35em plus 0.10em minus 0.08em}
\setlength{\emergencystretch}{3em}
\setlist{itemsep=0.25em, topsep=0.40em, parsep=0pt}
\allowdisplaybreaks[2]
% === 错题本排版优化结束 ===
\begin{document}
\section*{北京交通大学2025年数学分析真题}
\section*{6．北京交通大学2025年数学分析试卷}

\subsection{题目1}
1．设 \(f(x)\) 在 \(\mathbb{R}\) 上可导，若 \(f(x)\) 有两个零点，试证明在这两个零点之间必有 \(f(x)+f^{\prime}(x)\) 的零点．

\subsection{题目2}
2．设 \(0<x<1\) ，证明：\(e^{2 x}<\frac{x+1}{1-x}\) ．

\subsection{题目3}
3．已知 \(a_{0}=1, a_{n+1}=\frac{1}{a_{n}^{3}+4}, n=0,1,2, \cdots\) ．证明：

\begin{itemize}
  \item[（1）] \(\left\{a_{n}\right\}\) 收敛．
  \item[（2）] \(\left\{a_{n}\right\}\) 的极限是 \(x^{4}+4 x-1=0\) 的唯一正根．
\end{itemize}

\subsection{题目4}
4．求极限 \(\lim _{n \rightarrow \infty}\left[\left(1+\frac{1}{n}\right)\left(1+\frac{2}{n}\right) \cdots\left(1+\frac{n}{n}\right)\right]^{\frac{1}{n}}\) 。

\subsection{题目5}
5．已知 \(u=f(r)\) 二阶连续可微，\(r=\sqrt{x_{1}^{2}+x_{2}^{2}+\cdots+x_{n}^{2}}, r \neq 0, n \geq 1\) 。

\begin{itemize}
  \item[（1）] 试证明：\(\sum_{k=1}^{n} \frac{\partial^{2} u}{\partial x_{k}^{2}}=\frac{\mathrm{d}^{2} u}{\mathrm{~d} r^{2}}+\frac{n-1}{r} \frac{\mathrm{~d} u}{\mathrm{~d} r}(r \neq 0)\) ．
  \item[（2）] 求 \(\sum_{k=1}^{n} \frac{\partial^{2} u}{\partial x_{k}^{2}}=0\) 的解。
\end{itemize}

\subsection{题目6}
6．求最大的负数 \(A\) 和最小的正数 \(B\) ，使得不等式 \(\frac{A}{x y} \leq \ln (x+y) \leq \frac{B}{x^{2}+y^{2}}\) 对任意 \(x>0, y>0\) 恒成立。

\subsection{题目7}
7．求重积分 \(\iint_{D}|x-y+2| \mathrm{d} x \mathrm{~d} y\) ，其中 \(D: x^{2}+y^{2} \leq 4\) ．

\subsection{题目8}
8．求曲面积分

\[
\iint_{S}(x-y) \mathrm{d} y \mathrm{~d} z+\left(y^{2}-z^{2}\right) \mathrm{d} z \mathrm{~d} x+\left(z^{3}-x^{3}\right) \mathrm{d} x \mathrm{~d} y .
\]

其中 \(S\) 是锥面 \(\frac{x^{2}}{a^{2}}+\frac{y^{2}}{b^{2}}=z^{2}(a, b>0)\) 在 \(z=0\) 和 \(z=1\) 之间的部分，方向取下侧．

\subsection{题目9}
9．设 \(f_{0}(x)\) 在 \([a, b]\) 上连续，记 \(f_{n+1}(x)=\int_{a}^{x} f_{n}(t) \mathrm{d} t\) ，证明：函数项级数 \(\sum_{n=0}^{\infty} f_{n}(x)\) 在 \([a, b]\) 上一致收敛。

\subsection{题目10}
10．计算级数 \(\sum_{n=1}^{\infty} \frac{1}{2^{n}(2 n-1)}\) ．

\subsection{题目11}
11．设对任意的 \(n=1,2, \cdots, u_{n}(x)\) 在 \(x=c\) 处连续，但级数 \(\sum_{n=1}^{\infty} u_{n}(c)\) 发散，证明：对任意的 \(\delta>0\) ，函数项级数 \(\sum_{n=1}^{\infty} u_{n}(x)\) 在 \((c, c+\delta)\) 上非一致收敛。

\subsection{题目12}
12．设 \(f(x)\) 有连续导数，\(f_{n}(x)=n\left[f\left(x+\frac{1}{n}\right)-f(x)\right]\) ，证明：\(\left\{f_{n}(x)\right\}\) 在任意闭区间上均一致收敛．
\end{document}
